
\section{\extrahard\ \trGrassmannian{}：第一种构造}

\trGrassmannian{}是一种有用的几何构造，堪称「线性代数底下的那个几何对象」。在\trField{} $K=\mathbb{R}$ 或 $\mathbb{C}$ 上的（经典）几何中，正如\trProjectiveSpace{}参数化给定的 $n$ 维\trVectorSpace{}中的一维\trSubspace{}，\trGrassmannian{}参数化 $n$ 维空间中的 $k$ 维\trSubspace{}。\trGrassmannian{} $G(k,n)$ 是一个维数为 $k(n-k)$ 的\trManifold{}（在该\trField{}上）。\trManifold{}结构由如下方式给出。给定 $n$ 维空间的一组基 $(v_1,\dots,v_n)$，「大多数」$k$-平面（k-plane）都可以描述为由下面这 $k$ 个向量张成的\trSubspace{}：
\begin{equation*}
\left\langle v_1+\sum_{i=k+1}^{n}a_{1i}v_i,\;\;
v_2+\sum_{i=k+1}^{n}a_{2i}v_i,\;\;\dots,\;\;
v_k+\sum_{i=k+1}^{n}a_{ki}v_i\right\rangle.
\tag{7.7.0.1}\label{eq:grassmannian-basis-k-plane}
\end{equation*}
（你能描述哪些 $k$-平面\emph{不}属于这种形式吗？提示：\trRowReducedEchelonForm{}（row-reduced echelon form）。旁注：按\trNormalForm{}（normal form）对 $G(k,n)$ 作出的\trStratification{}（stratification），正是\trGrassmannian{}的\trSchubertCell{}（Schubert cells）分解。你或许能利用\trNormalForm{}证明：每个\trSchubertCell{}都\trIsomorphic{}于一个\trAffine{}空间。）这种形式的任何 $k$-平面都可以被唯一地如此描述。我们据此把这种形式的诸 $k$-平面等同于\trManifold{} $K^{k(n-k)}$（以诸 $a_{ji}$ 为\trCoordinate）。这是\trGrassmannian{}上的一大块\trAffine{}片（称作相对于这组基的「开\trSchubertCell{}」）。当诸 $v_i$ 变动时，这些片覆盖整个\trGrassmannian{}（为什么？），而且诸\trManifold{}结构彼此相合（这是一个更难的事实）。

我们现在于代数几何中、在一个\trRing{} $A$ 上来\emph{定义}\trGrassmannian。设 $v=(v_1,\dots,v_n)$ 是 $A^{\oplus n}$ 的一组基。更精确地说：诸 $v_i\in A^{\oplus n}$，并且由 $(a_1,\dots,a_n)\mapsto a_1v_1+\dots+a_nv_n$ 给出的映射 $A^{\oplus n}\to A^{\oplus n}$ 是一个\trIsomorphism。

\textbf{7.7.A.（练习）}\quad 证明：任何两组基都由 $A$ 上一个可逆的 $n\times n$ 矩阵相联系——所谓「可逆」，意即该矩阵元素取自 $A$，且其\trDeterminant{}是 $A$ 的一个可逆元。

对每一组这样的基 $v$，我们考虑\trScheme{} $U_v\cong\mathbb{A}^{k(n-k)}_A$，其\trCoordinate{}为诸 $a_{ji}$（$k+1\le i\le n$，$1\le j\le k$），我们把它想象成对应于由\eqref{eq:grassmannian-basis-k-plane}中诸向量张成的那一 $k$-平面。

\textbf{7.7.B.（练习）}\quad 给定两组基 $v$ 与 $w$，请说明如何沿适当的\trOpenSet{}把 $U_v$ 与 $U_w$ 粘合起来。你会发现下面这点或许带来方便：使用\trCoordinate{} $a_{ji}$（其 $i$ 从 $1$ 跑到 $n$，而不只是从 $k+1$ 跑到 $n$），但在 $i\le k$ 时强制约定 $a_{ji}=\delta_{ji}$（即 $i=j$ 时取 $1$，否则取 $0$）。这一约定与\trProjectiveSpace{}诸片上的\trCoordinate{} $x_{i/j}$ 相类似（\S4.4.9）。提示：$U_v$ 中相关的那个\trOpenSubset，将是某个\trDeterminant{}不消失之处。

\textbf{7.7.C.（练习/定义）}\quad 通过检验三重交，验证这些片（跑遍一切可能的基）能够粘合成单个\trScheme（\cref{ex:glue-arbitrary-schemes}）。这就是\trRing{} $A$ 上的\textbf{\trGrassmannian} $G(k,n)$。由于它可以解读为 $\mathbb{P}^{n-1}_A$ 中线性「$\mathbb{P}^{k-1}_A$」所成之空间，故也常写作 $G(k-1,n-1)$。（你将会看到这一记号确属明智，例如在\mbox{练习~12.2.N}中。）

尽管这个定义直白得令人愉快（\trGrassmannian{}被一片片 $\mathbb{A}^{k(n-k)}$ 覆盖是立见之事），或许也比我们在 \S4.4.9 中对\trProjectiveSpace{}的原始定义更「自然」（我们并不对基作选取；我们动用\emph{一切}基），它仍有数点不能令人满意。事实上，\trGrassmannian{}是\trProjective{}的；单看这个定义，这一点并不显然。此外，\trGrassmannian{}还自带一个到 $\mathbb{P}^{\binom{n}{k}-1}$ 的自然\trMorphism（它是\trClosedEmbedding{}的一个例子；这一观念我们将在定义~9.1.1 中遇到），称为\trPluckerEmbedding{}（Plücker embedding）。最后，$\mathrm{GL}_n$ 在「$n$ 维空间中的 $k$-平面」所成之空间上有一个作用，因此我们理应能在我们的代数化身中看到这个作用。我们将在 \S16.4 中处理这些问题，办法是给出一份更好的描述：把它作为\trModuliSpace{}。

\textbf{7.7.1.}\ \emph{（部分）\trFlagVariety{}。}\quad 正如\trGrassmannian{}「参数化」$n$ 维空间中的 $k$-平面，\textbf{\trFlagVariety{}}（flag variety）参数化「旗」（flags）：$n$ 维空间中\trSubspace{}的嵌套序列
\[
F_0\subset F_1\subset\dots\subset F_n,
\]
其中 $\dim F_i=i$。把这两者一并推广的，是与数据 $0\le a_1<\dots<a_\ell\le n$ 相关联的\textbf{\trPartialFlagVariety{}}（partial flag variety）这一观念，它参数化 $n$ 维空间中\trSubspace{}的嵌套序列
\[
F_{a_1}\subset\dots\subset F_{a_\ell},
\]
其中 $\dim F_{a_i}=a_i$。\S7.7 中的全部讨论，你都应当能够推广到这个场合。
